% 一条原弧的两条残量弧：容量 3、流量 2 -> 正向残量 1，反向残量 2。
\begin{tikzpicture}[>=Stealth,
  v/.style={circle,draw=FigureStroke,fill=FigureBlueFill,line width=.8pt,minimum size=6.5mm,inner sep=0pt},
  w/.style={fill=white,inner sep=.8pt,font=\figurefont\fontsize{7.5}{9}\selectfont},
  lab/.style={font=\figurefont\fontsize{7.5}{9}\selectfont}]
% Panel labels share one baseline, 3 mm below the full row of content.
\begin{scope}[local bounding box=panel-content]
\begin{scope}
  \node[v] (u) at (0,0) {$u$}; \node[v] (v) at (2.4,0) {$v$};
  \draw[->,draw=FigureStroke,line width=.8pt] (u)--node[w]{$2/3$}(v);
  \coordinate[overlay] (panel-a-center) at (1.2,-1.05);
\end{scope}
\begin{scope}[xshift=4.2cm]
  \node[v] (u) at (0,0) {$u$}; \node[v] (v) at (2.4,0) {$v$};
  \draw[->,draw=FigureBlue,line width=1.2pt] (u) to[bend left=28] node[w]{$1$}(v);
  \draw[->,draw=FigureGreen,dashed,line width=1.2pt] (v) to[bend left=28] node[w]{$2$}(u);
  \node[lab,text=FigureBlue] at (1.2,.85) {正向：还能追加};
  \node[lab,text=FigureGreen] at (1.2,-.85) {反向：能撤回};
  \coordinate[overlay] (panel-b-center) at (1.2,-1.4);
\end{scope}
\end{scope}
\node[font=\figurefont\fontsize{8}{10}\selectfont,text=FigureInk,anchor=north,inner sep=0pt,outer sep=0pt] at ([yshift=-3mm]panel-a-center |- panel-content.south) {(a)};
\node[font=\figurefont\fontsize{8}{10}\selectfont,text=FigureInk,anchor=north,inner sep=0pt,outer sep=0pt] at ([yshift=-3mm]panel-b-center |- panel-content.south) {(b)};
\end{tikzpicture}
